\chapter{Conclusão}
\label{cap:conclusao}

Este guia percorreu a cadeia completa de uma análise causal quantitativa:
da pergunta de política ao número numa tabela. Cada escolha metodológica ao longo desse caminho reflete uma hipótese sobre o mundo, e cada hipótese carrega um custo que precisa ser explicitado e defendido.

O argumento central pode ser enunciado de forma direta: \textbf{causalidade não é uma propriedade dos dados, mas das hipóteses que fazemos sobre como os dados foram gerados}. Dados observacionais não ``falam por si''; eles precisam ser lidos à luz de um argumento de identificação. Esse argumento determina o que é estimado, e a estimação determina com que precisão. Sem o primeiro elo, os outros dois não produzem nada interpretável.


\section{A Cadeia de Análise: Um Balanço}

\subsection*{Identificação}

O Capítulo~\ref{cap:identificacao} mostrou que o problema fundamental da
inferência causal --- a impossibilidade de observar o contrafactual --- não
pode ser eliminado, apenas contornado. O contorno exige hipóteses sobre o
processo de seleção ao tratamento. A exogeneidade incondicional (satisfeita
por randomização) elimina o viés de seleção por construção. A independência
condicional (CIA) exige que todas as variáveis que simultaneamente afetam
tratamento e resultado estejam disponíveis nos dados --- uma hipótese
substantiva, não testável diretamente.

A lição prática: antes de estimar qualquer coisa, o pesquisador precisa
responder à pergunta ``por que comparar esses dois grupos é informativo sobre
o efeito causal?'' A resposta a essa pergunta é o argumento de identificação.
Sem ele, nenhum refinamento econométrico subsequente produz interpretação
causal.

\subsection*{Estimação}

O Capítulo~\ref{cap:estimacao} mostrou que identificação e estimação são
problemas distintos. Mesmo quando um estimando é identificado, o estimador
precisa efetivamente calculá-lo. A forma funcional importa: OLS com controles
lineares pode calcular uma média ponderada adversa dos efeitos locais quando
os pesos implícitos diferem dos pesos de interesse. O IPW oferece robustez à
forma funcional de $\mathbb{E}[Y \mid D, X]$ ao custo de depender do
propensity score. Estimadores duplamente robustos combinam os dois modelos
e são consistentes se pelo menos um estiver correto.

\subsection*{Inferência}

O Capítulo~\ref{cap:inferencia} mostrou que a escolha do estimador de
variância é tão substantiva quanto a escolha do estimador do parâmetro.
Erros-padrão clássicos supõem independência e homoscedasticidade ---
hipóteses raramente plausíveis em dados de políticas públicas. O clustering
é obrigatório quando o tratamento é atribuído em grupo. Com poucos clusters,
o wild cluster bootstrap é a solução preferida.

O p-valor mede uma coisa específica: a probabilidade de observar dados tão
extremos quanto os observados, se a hipótese nula fosse verdadeira. Ele não
mede o tamanho do efeito e não substitui o intervalo de confiança. Reportar
o tamanho do efeito com seu intervalo de confiança é sempre mais informativo
do que reportar apenas o p-valor.


\section{O Experimento STAR como Laboratório}

O Capítulo~\ref{cap:star} demonstrou todos esses princípios em um único
exemplo empírico. O STAR mostra que:

\begin{itemize}[leftmargin=*]
  \item A randomização resolve o problema de identificação dentro dos blocos
    experimentais --- mas não elimina a necessidade de explicitar as hipóteses
    (SUTVA pode ser questionado; a qualidade da implementação importa).
  \item A especificação com efeitos fixos de escola alinha a estimação ao
    design experimental, melhorando precisão sem comprometer validade.
  \item A escolha entre erros clássicos, HC e clusterizados produz diferenças
    substanciais nos erros-padrão --- diferenças com consequências reais para
    as conclusões.
  \item Testar quatro disciplinas sem correção para múltiplos testes infla a
    taxa de falsos positivos. Os ajustes de Bonferroni e BH-FDR revelam quais
    resultados são robustos.
  \item A inferência por permutação usa diretamente o processo de aleatorização
    e produz p-valores exatos em amostras finitas.
  \item Os testes de falsificação fornecem evidência circunstancial sobre a
    qualidade do desenho, mas não provam identificação válida.
\end{itemize}


\section{O Que Este Guia Não Cobre}

Este guia apresentou os fundamentos da inferência causal em contextos
experimentais e observacionais com seleção em observáveis. Os casos em que
a CIA não pode ser mantida exigem estratégias quasi-experimentais que exploram
variação exógena de outras naturezas. Esses métodos são o objeto dos demais
guias da Série Avaliação na Prática:

\begin{itemize}[leftmargin=*]
  \item \textbf{Regressão com Descontinuidade (RDD)}: explora pontos de corte
    em regras de elegibilidade para comparações localmente aleatórias.
  \item \textbf{Diferenças em Diferenças (DiD)}: utiliza variação temporal no
    tratamento para controlar por diferenças fixas não observadas entre grupos.
  \item \textbf{Pareamento (\textit{Matching})}: constrói um grupo de controle
    sintético com covariadas similares às do grupo tratado.
  \item \textbf{Controle Sintético (SCM)}: combina unidades de controle para
    replicar a trajetória pré-tratamento da unidade tratada.
  \item \textbf{Variáveis Instrumentais (IV)}: utiliza uma variável que afeta
    o tratamento mas não afeta diretamente o resultado.
\end{itemize}

\begin{conexaobox}[Série Avaliação na Prática]
Os guias de RDD, DiD (dois volumes), Pareamento, SCM, Variáveis Instrumentais,
Poder Estatístico e Fontes de Dados estão disponíveis na Série Avaliação na
Prática do FGV CLEAR. Cada guia pode ser lido de forma independente, mas os
conceitos de identificação, estimação e inferência desenvolvidos aqui são
pressupostos em todos eles.
\end{conexaobox}


\section{Uma Nota Final}

A econometria causal não é uma caixa de ferramentas que, aplicada
corretamente, produz respostas definitivas. É um conjunto de disciplinas
intelectuais para raciocinar com rigor sobre o que podemos e o que não podemos
aprender a partir de dados. Mesmo o experimento mais bem controlado carrega
hipóteses que precisam ser defendidas, não apenas assumidas.

O hábito que este guia procura cultivar é uma pergunta, repetida a cada estimativa: \textit{o que exatamente estou estimando? Sob que hipóteses esse número tem interpretação causal? Quão plausíveis são essas hipóteses neste contexto?} Um pesquisador que faz essas perguntas sistematicamente produz avaliações mais informativas do que aquele que domina ferramentas sofisticadas sem questionar o que elas calculam.

Quando essas perguntas são levadas a sério --- e as hipóteses são explicitadas e debatidas em vez de varridas para baixo do tapete ---, a inferência causal oferece a melhor base disponível para decisões de política pública informadas por evidência.